\section{Discussion and Future Work}
\label{sec:lim-fut}

In this section, we briefly discuss limitations and future directions.

\paragraph{Looped Architectures.} While several design choices around looped architectures have been guided by small-scale empirical results, a deep investigation of loop-unit placement \cite{jacobs2026blockrecurrentdynamicsvisiontransformers}, composition (e.g., number of parameters in the recurrent unit and usage of different architectures), and extreme looping (e.g., increasing mean recurrence to deeper depths) at a larger scale is warranted. Within our dynamical systems framework, the use of different discretizations, full-rank parameterizations, and recurrent update rules warrants investigation to enable recurrence at larger depths.

\paragraph{Scaling.}
While we find Parcae to induce predictable, optimal scaling laws for layer looping, our observations are limited to small architectures. It remains to be seen if Parcae compares favorably when scaling these observations to large FLOP budgets and parameterizations. We are also interested in the interplay of parameters, data, and recurrence as orthogonal axes, and how they should be efficiently scaled together. Finally, one limitation of looping is that, as \meanrecurrence increases, the number of test-time steps required to achieve equivalent quality increases. An investigation of techniques that maintain quality with fewer inference time steps is an interesting future direction.







\section{Conclusion}
\label{sec:conclusion}

In this work, we study the stability of looped models through a dynamical systems framework and propose \textbf{\sysname}, a stable looped architecture that prevents residual explosion by constraining the spectral norm of the injection parameters.
\sysname outperforms data- and parameter-matched prior looped models and baseline Transformers, matching downstream quality of models up to twice its size.
We further establish scaling laws for looping: FLOP-optimal training increases looping and data in tandem following predictable power laws, while test-time looping follows a saturating exponential decay law, yielding a unified scaling law connecting training and inference compute. 

